\documentclass[letterpaper]{article}

\usepackage[letterpaper,margin=1in]{geometry}

% Compile from the repository root so this path resolves to the pinned WSP
% submodule.
\input{wsp/testing/test-case-library.tex}

\begin{document}
\def\TCID{TC-0002}
\def\TCTitle{Keyboard and INKEY\$}
\def\TCPurpose{Verify the strings \texttt{INKEY\$} returns, the BIOS codes of
the extended keys with Shift, Ctrl and Alt, the translation of Windows key
presses and the 15-key type-ahead buffer.}
\def\TCPriority{\TCPriorityHigh}

\def\TCRequirementRef{REQ-0002}
\def\TCDesignRef{Not applicable.}
\def\TCFeatureRef{\texttt{include/wfc/brun71/keyboard.hpp};
\texttt{src/keyboard.cpp}; \texttt{tests/keyboard\_test.cpp}.}
\def\TCTestTechnique{Decision table of key and modifier combinations;
boundary values of the buffer (empty, 15 keys, one more).}

\def\TCPreconditions{The library is built and the test program
\texttt{brun71\_keyboard\_test} exists.}
\def\TCEnvironment{Any supported target; the translation is plain code with
no console. CTest runs the program with a 60 second timeout.}
\def\TCAssumptions{The expected codes are the PC BIOS keyboard codes.}

\def\TCInputData{Key presses (virtual-key code, character, Shift, Ctrl, Alt)
for letters, control keys, cursor keys, function keys, AltGr characters and
modifier keys alone.}
\def\TCInitialState{An empty key buffer.}

\def\TCProcedure{\begin{enumerate}
  \item Run \texttt{brun71\_keyboard\_test} (CTest
    \texttt{TC-0002-keyboard}, in WFC \texttt{wfc-brun71-TC-0002-keyboard}).
  \item Check the \texttt{INKEY\$} string of a character and of an extended
    key.
  \item Check the translation of characters, Ctrl+letter, AltGr, a character
    outside code page 437 and the modifier keys alone.
  \item Check the codes of the cursor, editing and function keys alone and
    with Shift, Ctrl and Alt; Shift+Tab; Alt with letters, digits, minus and
    equals.
  \item Check that the buffer keeps 15 keys in order, refuses a 16th and
    returns the empty string when empty.
\end{enumerate}}
\def\TCExpectedResult{The program exits with status 0 and prints
\texttt{keyboard: ok}; no \texttt{FAIL} line is written.}
\def\TCPostConditions{None; the test writes no files.}

\TCTemplate

\end{document}
